\documentclass[12pt,a4paper]{article}
\usepackage{../../../myconfig}

\begin{document}

% =========================================================================
% BÀI TẬP VỀ NHÀ: PHƯƠNG PHÁP TỌA ĐỘ HÓA / HHKG (20 CÂU MỚI 100%)
% =========================================================================
\titlebox{Chuyên đề: Oxyz}{Phương pháp tọa độ hóa}{Oxyz: Phương pháp tọa độ hóa}

\vspace{-0.25cm}

% -------------------------------------------------------------------------
% PHẦN 1: KHOẢNG CÁCH GIỮA HAI ĐƯỜNG THẲNG CHÉO NHAU
% -------------------------------------------------------------------------

% Câu 1 (Câu 5 - Dạng 2 - Chuyên đề 4)
\questionLinesManual{0}{1}{%
Cho hình chóp $S.ABCD$ có đáy là hình chữ nhật, $AB = a$, $BC = 2a$, $SA$ vuông góc với mặt phẳng đáy và $SA = a$. Khoảng cách giữa hai đường thẳng $AC$ và $SB$ bằng
\begin{tasks}(4)
    \task $\dfrac{\sqrt{6}a}{2}$.
    \task $\dfrac{2a}{3}$.
    \task $\dfrac{a}{2}$.
    \task $\dfrac{a}{3}$.
\end{tasks}
}{}

\vspace{0.15cm}

% Câu 2 (Câu 6 - Dạng 2 - Chuyên đề 4)
\questionLinesManual{0}{2}{%
Cho tứ diện $OABC$ có $OA, OB, OC$ đôi một vuông góc với nhau và $OA = OB = a$, $OC = 2a$. Gọi $M$ là trung điểm của $AB$. Khoảng cách giữa hai đường thẳng $OM$ và $AC$ bằng
\begin{tasks}(4)
    \task $\dfrac{2\sqrt{5}a}{5}$.
    \task $\dfrac{\sqrt{2}a}{2}$.
    \task $\dfrac{2a}{3}$.
    \task $\dfrac{\sqrt{2}a}{3}$.
\end{tasks}
}{}

\vspace{0.15cm}

% Câu 3 (Câu 10 - Dạng 2 - Chuyên đề 4)
\questionLinesManual{0}{3}{%
Cho tứ diện $OABC$ có $OA, OB, OC$ đôi một vuông góc với nhau, $OA = a$ và $OB = OC = 2a$. Gọi $M$ là trung điểm của $BC$. Khoảng cách giữa hai đường thẳng $OM$ và $AB$ bằng
\begin{tasks}(4)
    \task $\dfrac{\sqrt{6}a}{3}$.
    \task $a$.
    \task $\dfrac{2\sqrt{5}a}{5}$.
    \task $\dfrac{\sqrt{2}a}{2}$.
\end{tasks}
}{}

\vspace{0.15cm}

% Câu 4 (Câu 14 - Dạng 2 - Chuyên đề 4)
\questionLinesManual{0}{4}{%
Cho hình chóp $S.ABC$ có tam giác $ABC$ vuông tại $B$, $\widehat{C} = 60^\circ$, $AC = 2$, $SA \perp (ABC)$ và $SA = 1$. Gọi $M$ là trung điểm của $AB$. Khoảng cách $d$ giữa hai đường thẳng $SM$ và $BC$ bằng
\begin{tasks}(4)
    \task $d = \dfrac{\sqrt{21}}{7}$.
    \task $d = \dfrac{2\sqrt{21}}{7}$.
    \task $d = \dfrac{\sqrt{21}}{3}$.
    \task $d = \dfrac{2\sqrt{21}}{3}$.
\end{tasks}
}{}

\vspace{0.15cm}

% Câu 5 (Câu 20 - Dạng 2 - Chuyên đề 4)
\questionLinesManual{0}{5}{%
Cho hình chóp $S.ABC$ có đáy $ABC$ là tam giác vuông cân tại $A$, $AB = a$. Cạnh bên $SA$ vuông góc với mặt phẳng đáy và $SA = a$. Gọi $M$ là trung điểm của $BC$. Khoảng cách giữa hai đường thẳng $AC$ và $SM$ bằng
\begin{tasks}(4)
    \task $\dfrac{\sqrt{3}a}{3}$.
    \task $\dfrac{\sqrt{2}a}{2}$.
    \task $\dfrac{a}{2}$.
    \task $\dfrac{\sqrt{5}a}{5}$.
\end{tasks}
}{}

\vspace{0.15cm}

% -------------------------------------------------------------------------
% PHẦN 2: KHOẢNG CÁCH TỪ ĐIỂM ĐẾN MẶT PHẲNG
% -------------------------------------------------------------------------

% Câu 6 (Câu 6 - Dạng 1 - Chuyên đề 4)
\questionLinesManual{0}{6}{%
Cho hình chóp $S.ABC$ có đáy là tam giác vuông đỉnh $B$, $AB = a$, $SA$ vuông góc với mặt phẳng đáy và $SA = a$. Khoảng cách từ điểm $A$ đến mặt phẳng $(SBC)$ bằng
\begin{tasks}(4)
    \task $\dfrac{a\sqrt{6}}{3}$.
    \task $\dfrac{a\sqrt{2}}{2}$.
    \task $\dfrac{a}{2}$.
    \task $a$.
\end{tasks}
}{}

\vspace{0.15cm}

% Câu 7 (Câu 11 - Dạng 1 - Chuyên đề 4)
\questionLinesManual{0}{7}{%
Cho hình chóp $S.ABCD$ có đáy là hình vuông cạnh $\sqrt{3}a$, $SA$ vuông góc với mặt phẳng đáy và $SA = a$. Khoảng cách từ điểm $A$ đến mặt phẳng $(SBC)$ bằng
\begin{tasks}(4)
    \task $\dfrac{\sqrt{6}a}{6}$.
    \task $\dfrac{\sqrt{3}a}{3}$.
    \task $\dfrac{\sqrt{5}a}{3}$.
    \task $\dfrac{\sqrt{3}a}{2}$.
\end{tasks}
}{}

\vspace{0.15cm}

% Câu 8 (Câu 13 - Dạng 1 - Chuyên đề 4)
\questionLinesManual{0}{8}{%
Cho hình chóp $S.ABCD$ có $SA \perp (ABCD)$, đáy $ABCD$ là hình chữ nhật. Biết $AD = 2a$, $SA = a$. Khoảng cách từ điểm $A$ đến mặt phẳng $(SCD)$ bằng
\begin{tasks}(4)
    \task $\dfrac{3a}{\sqrt{7}}$.
    \task $\dfrac{3a\sqrt{2}}{2}$.
    \task $\dfrac{2a}{\sqrt{5}}$.
    \task $\dfrac{2a\sqrt{3}}{3}$.
\end{tasks}
}{}

\vspace{0.15cm}

% Câu 9 (Câu 15 - Dạng 1 - Chuyên đề 4)
\questionLinesManual{0}{9}{%
Cho hình chóp tứ giác đều $S.ABCD$ có cạnh đáy bằng $a$ và chiều cao bằng $a\sqrt{2}$. Tính khoảng cách $d$ từ tâm $O$ của đáy $ABCD$ đến một mặt bên theo $a$.
\begin{tasks}(4)
    \task $d = \dfrac{2a\sqrt{5}}{3}$.
    \task $d = \dfrac{a\sqrt{3}}{2}$.
    \task $d = \dfrac{a\sqrt{5}}{2}$.
    \task $d = \dfrac{a\sqrt{2}}{3}$.
\end{tasks}
}{}

\vspace{0.15cm}

% Câu 10 (Câu 28 - Dạng 1 - Chuyên đề 4)
\questionLinesManual{0}{10}{%
Cho hình chóp $S.ABCD$ có đáy là hình vuông, $AB = a$, $SA$ vuông góc với mặt phẳng đáy và $SA = 2a$. Gọi $M$ là trung điểm của cạnh $CD$, khoảng cách giữa điểm $M$ và mặt phẳng $(SBD)$ bằng
\begin{tasks}(4)
    \task $\dfrac{2a}{3}$.
    \task $\dfrac{a}{\sqrt{2}}$.
    \task $\dfrac{a}{2}$.
    \task $\dfrac{a}{3}$.
\end{tasks}
}{}

\vspace{0.15cm}

% -------------------------------------------------------------------------
% PHẦN 3: GÓC GIỮA HAI MẶT PHẲNG
% -------------------------------------------------------------------------

% Câu 11 (Câu 4 - Dạng 3 - Chuyên đề 3)
\questionLinesManual{0}{11}{%
Cho hình chóp $S.ABCD$ có đáy $ABCD$ là hình chữ nhật, $AB = a$, $AD = SA = 2a$, $SA \perp (ABCD)$. Tính tang của góc giữa hai mặt phẳng $(SBD)$ và $(ABCD)$.
\begin{tasks}(4)
    \task $\dfrac{\sqrt{5}}{2}$.
    \task $\sqrt{5}$.
    \task $\dfrac{1}{\sqrt{5}}$.
    \task $\dfrac{2}{\sqrt{5}}$.
\end{tasks}
}{}

\vspace{0.15cm}

% Câu 12 (Câu 5 - Dạng 3 - Chuyên đề 3)
\questionLinesManual{0}{12}{%
Cho hình chóp $S.ABCD$ có đáy $ABCD$ là hình thoi tâm $O$, đường thẳng $SO$ vuông góc với mặt phẳng $(ABCD)$. Biết $AB = SB = a$, $SO = \dfrac{a\sqrt{6}}{3}$. Tìm số đo của góc giữa hai mặt phẳng $(SAB)$ và $(SAD)$.
\begin{tasks}(4)
    \task $30^\circ$.
    \task $45^\circ$.
    \task $60^\circ$.
    \task $90^\circ$.
\end{tasks}
}{}

\vspace{0.15cm}

% Câu 13 (Câu 11 - Dạng 3 - Chuyên đề 3)
\questionLinesManual{0}{13}{%
Cho hình chóp $S.ABCD$ có đáy $ABCD$ là hình thoi cạnh $a$, $BD = a$. Cạnh bên $SA$ vuông góc với mặt đáy và $SA = \dfrac{a\sqrt{6}}{2}$. Tính góc giữa hai mặt phẳng $(SBC)$ và $(SCD)$.
\begin{tasks}(4)
    \task $60^\circ$.
    \task $120^\circ$.
    \task $45^\circ$.
    \task $90^\circ$.
\end{tasks}
}{}

\vspace{0.15cm}

% Câu 14 (Câu 27 - Dạng 3 - Chuyên đề 3)
\questionLinesManual{0}{14}{%
Cho hình lập phương $ABCD.A'B'C'D'$. Côsin góc giữa hai mặt phẳng $(A'BC)$ và $(ABC')$ bằng
\begin{tasks}(4)
    \task $\dfrac{\sqrt{3}}{2}$.
    \task $\dfrac{\sqrt{2}}{2}$.
    \task $0$.
    \task $\dfrac{1}{2}$.
\end{tasks}
}{}

\vspace{0.15cm}

% Câu 15 (Câu 23 - Dạng 3 - Chuyên đề 3)
\questionLinesManual{0}{15}{%
Cho hình chóp $S.ABC$ có đáy $ABC$ là tam giác đều cạnh $a$, $SA$ vuông góc với mặt phẳng $(ABC)$ và $SA = \dfrac{a}{2}$. Góc giữa mặt phẳng $(SBC)$ và mặt phẳng $(ABC)$ bằng
\begin{tasks}(4)
    \task $45^\circ$.
    \task $90^\circ$.
    \task $30^\circ$.
    \task $60^\circ$.
\end{tasks}
}{}

\vspace{0.15cm}

% -------------------------------------------------------------------------
% PHẦN 4: GÓC GIỮA HAI ĐƯỜNG THẲNG / ĐƯỜNG THẲNG VÀ MẶT PHẲNG
% -------------------------------------------------------------------------

% Câu 16 (Câu 5 - Dạng 1 - Chuyên đề 3)
\questionLinesManual{0}{16}{%
Cho hình lập phương $ABCD.A'B'C'D'$. Gọi $M$ là trung điểm của $B'C'$. Góc giữa hai đường thẳng $AM$ và $BC'$ bằng
\begin{tasks}(4)
    \task $45^\circ$.
    \task $90^\circ$.
    \task $30^\circ$.
    \task $60^\circ$.
\end{tasks}
}{}

\vspace{0.15cm}

% Câu 17 (Câu 9 - Dạng 1 - Chuyên đề 3)
\questionLinesManual{0}{17}{%
Cho hình lập phương $ABCD.A'B'C'D'$. Gọi $M, N$ lần lượt là trung điểm của $AD, BB'$. Côsin của góc hợp bởi $MN$ và $AC'$ bằng
\begin{tasks}(4)
    \task $\dfrac{\sqrt{3}}{3}$.
    \task $\dfrac{\sqrt{2}}{3}$.
    \task $\dfrac{\sqrt{5}}{3}$.
    \task $\dfrac{\sqrt{2}}{4}$.
\end{tasks}
}{}

\vspace{0.15cm}

% Câu 18 (Câu 12 - Dạng 1 - Chuyên đề 3)
\questionLinesManual{0}{18}{%
Cho hình chóp đều $S.ABCD$ có tất cả các cạnh đều bằng $a$. Gọi $M, N$ lần lượt là trung điểm của $AD$ và $SD$. Số đo của góc giữa hai đường thẳng $MN$ và $SC$ là
\begin{tasks}(4)
    \task $45^\circ$.
    \task $60^\circ$.
    \task $30^\circ$.
    \task $90^\circ$.
\end{tasks}
}{}

\vspace{0.15cm}

% Câu 19 (Câu 1 - Dạng 2 - Chuyên đề 3)
\questionLinesManual{0}{19}{%
Cho hình chóp $S.ABCD$ có đáy là hình vuông cạnh $\sqrt{3}a$, $SA$ vuông góc với mặt phẳng đáy và $SA = \sqrt{2}a$. Góc giữa đường thẳng $SC$ và mặt phẳng $(ABCD)$ bằng
\begin{tasks}(4)
    \task $45^\circ$.
    \task $60^\circ$.
    \task $30^\circ$.
    \task $90^\circ$.
\end{tasks}
}{}

\vspace{0.15cm}

% Câu 20 (NÂNG CAO - ẨN CHIỀU CAO h: Câu 21 - Dạng 3 - Chuyên đề 3)
\questionLinesManual{0}{20}{%
Cho hình chóp $S.ABCD$ có đáy $ABCD$ là hình thoi cạnh $a$, $\widehat{ABC} = 120^\circ$. Cạnh bên $SA$ vuông góc với mặt phẳng đáy $(ABCD)$ và có độ dài chưa biết $SA = h$ ($h > 0$). Biết góc giữa hai mặt phẳng $(SBC)$ và $(SCD)$ bằng $60^\circ$. Chiều cao $h = SA$ của khối chóp bằng
\begin{tasks}(4)
    \task $SA = \dfrac{a\sqrt{6}}{4}$.
    \task $SA = a\sqrt{6}$.
    \task $SA = \dfrac{a\sqrt{6}}{2}$.
    \task $SA = \dfrac{a\sqrt{3}}{2}$.
\end{tasks}
}{}


% =========================================================================
% ĐÁP ÁN BÀI TẬP VỀ NHÀ (20 CÂU KHÔNG TRÙNG ĐIỂM DANH)
% =========================================================================
% 01. B | 06. B | 11. B | 16. A
% 02. C | 07. D | 12. D | 17. B
% 03. A | 08. C | 13. D | 18. D
% 04. A | 09. D | 14. D | 19. C
% 05. D | 10. D | 15. C | 20. A

\end{document}